\begin{figure}[p]
\centering
\includegraphics[width=\textwidth]{supplementary/curated/shunt_sensitivity.pdf}
\caption{\textbf{Dose, cable parameters and input mapping determine focal-shunt selectivity.} \textbf{A}, Shunt-minus-current-injection localization across fixed absolute doses at three membrane resistances ($R_m$ of 300, 1,000 and 15,000~$\Omega\,\mathrm{cm}^2$: one lightness step, marker and dash pattern each) and three background-conductance multipliers (column groups), on a symmetric-log axis that is linear within $\pm0.001$. The contrast is close to proportional to dose at every setting, and the $R_m$ ordering reverses with background: with no background the $R_m=15{,}000$ contrast stays near zero (slightly negative at 0.05 and 0.5~nS, an interval spanning zero at 5~nS), whereas at background multiplier four it is the largest of the three at every dose, with 8/8 cells positive. \textbf{B}, Shunt and matched-injection localization (grey lines pair cells) and the true relation against a reassigned foreign template (unpaired); positive favors the shunt or the true relation. \textbf{C}, Transport selectivity $S_k$ against full synaptic-gradient localization at unit input-conductance-normalized dose, 101 sites; descriptive site--regime points, $R_m$ colours as in \textbf{A}, background multipliers 0, 1 and 4 as filled, open and plus markers; dotted line, $S_k=1$. \textbf{D}, Signed census of the same sites at $R_m=1{,}000\,\Omega\,\mathrm{cm}^2$, background multiplier one, unit dose: fraction of descendant gradients attenuated, enhanced or sign reversed by shunt (attenuates only) and matched injection (enhances only), as in all 81 passive conditions. \textbf{E}, All mapped versus directly typed contacts (grey lines pair cells); the contrast increases in 8/8 cells. \textbf{F}, E/I scale and inhibitory-reversal sensitivity by row; open marker and band, the reference condition (E/I 0.35, reversal $-0.2$) and its interval; right column, cells with a positive contrast. \textbf{A,B,E,F}, means and 95\% cell-bootstrap intervals, eight cells.}
\label{fig:si_shunt_sensitivity}
\label{fig:supp_focal_matrix}
\label{fig:supp_focal_detail}
\end{figure}

\begin{figure}[p]
\centering
\includegraphics[width=\textwidth]{supplementary/curated/shunt_replication.pdf}
\caption{\textbf{Focal-shunt controls replicate structurally and under weak-channel linearization.} \textbf{A}, Focal shunt versus baseline-current-matched injection in 45 disjoint same-mouse cells (235 sites). \textbf{B}, True versus reassigned descendant relation templates in forty cells (230 sites); the true-relation column repeats the focal-shunt values of \textbf{A} for the forty cells that have a reassigned-relation partner. \textbf{C}, Direct E/I label coverage and the number of selected focal sites per cell, capped at 16 (one L5ET cell had 47 eligible sites); cell classes are ordinal colours within this panel. \textbf{D,E}, Separate initial eight-cell weak-channel ensemble: localization versus dose (dose divided by local input conductance) and paired shunt-minus-current contrasts. A fixed steady-state Jacobian (solid) is compared with the passive model (dashed); the two means differ by at most 0.010 localization units for the shunt and 0.023 for current injection at every dose, so the passive curve lies under the active one. Shunting attenuates descendants whereas matched current injection enhances them. Sixty-four channel draws and sites are averaged within each cell; intervals bootstrap cells and are narrower than the marker except at dose 4, and the contrast is positive in 8/8 cells at every dose (Wilcoxon $p=0.0078$). Weak conductances test local linearization near the passive response and do not establish robustness to regenerative dynamics or channel kinetics.}
\label{fig:si_shunt_replication}
\label{fig:supp_weak_channel_linearization}
\end{figure}
